\chapter{Fundamentação Teórica}
\label{chap:fundamentacao}

\epigraph{\itshape ``Não basta fazer a coisa certa da primeira vez. É preciso poder
fazê-la da mesma forma na milésima vez.''}{\normalfont --- princípio da engenharia de qualidade}

Este capítulo apresenta os fundamentos teóricos necessários para compreender e
contextualizar o framework proposto. A narrativa segue uma progressão natural: da
física do robô às abstrações de software que a encapsulam, e destas às métricas
que permitem avaliar se o sistema se comporta como projetado.

\section{Cinemática de Robôs com Acionamento Diferencial}
\label{sec:cinematica}

O TurtleBot4 Standard, plataforma de referência desta dissertação, opera com acionamento
diferencial: duas rodas traseiras independentes e uma roda dianteira livre (caster)
que segue passivamente o movimento. Nesse modelo, a velocidade linear $v$ e angular
$\omega$ do robô relacionam-se com as velocidades individuais das rodas $v_R$ e $v_L$
pela cinemática inversa \cite{siegwart2011}:

\begin{equation}
v = \frac{v_R + v_L}{2}, \quad \omega = \frac{v_R - v_L}{L}
\label{eq:cinematica_diferencial}
\end{equation}

\noindent onde $L$ é a distância entre as rodas (wheelbase). A integração dessas
velocidades ao longo do tempo produz a estimativa de pose $(x, y, \theta)$ --- o
que é denominado \textit{dead reckoning} ou odometria de roda.

\subsection{Acúmulo de Erro de Odometria}
\label{sec:erro_odometria}

A odometria por integração é inerentemente propensa a erros acumulados
\cite{thrun2005, siegwart2011}. Três fontes dominam:

\begin{enumerate}
    \item \textbf{Resolução dos encoders}: encoders incrementais possuem resolução
    finita; a discretização introduz erro por passo que, integrado ao longo de
    centenas de metros, pode exceder um metro de deriva posicional.

    \item \textbf{Slip de roda}: deslizamento entre roda e superfície viola a
    suposição de rolamento puro do modelo cinemático, introduzindo erro sistemático
    que varia com velocidade, carga e tipo de superfície.

    \item \textbf{Desalinhamento de rodas}: pequenas diferenças no diâmetro efetivo
    de cada roda introduzem curvatura sistemática mesmo em trajetos nominalmente retos.
\end{enumerate}

Para robôs de acionamento diferencial em superfície regular, erros de odometria pura
da ordem de 5--10\% do percurso total são típicos \cite{siegwart2011}. Percursos de
5\,m podem acumular até 50\,cm de erro de posição sem correção externa --- tornando o SLAM
indispensável para experimentos que exijam repetibilidade sub-decimétrica.

\section{Localização e Mapeamento Simultâneos (SLAM)}
\label{sec:slam}

O SLAM resolve o problema de localizar um robô em um mapa que está sendo construído
ao mesmo tempo \cite{durrantwhyte2006, thrun2005}. Formalmente, dado um histórico de
observações $z_{1:t}$ (scans LiDAR) e controles $u_{1:t}$ (odometria), o SLAM
estima a distribuição de probabilidade conjunta sobre a pose do robô $x_t$ e o mapa
$m$:

\begin{equation}
p(x_t, m \mid z_{1:t}, u_{1:t})
\label{eq:slam_probabilistico}
\end{equation}

Diferentes algoritmos aproximam essa distribuição de formas distintas. Filtros de
partículas (como o GMapping) representam a incerteza com um conjunto de hipóteses
amostradas; abordagens baseadas em grafos (como o SLAM Toolbox
\cite{macenski2021slam}) representam a trajetória do robô como um grafo de poses
e minimizam a energia do grafo via otimização esparsa. Esta última abordagem escala
melhor para ambientes grandes e é mais adequada para mapeamento lifelong.

\subsection{A Cadeia TF no ROS~2}
\label{sec:cadeia_tf}

No ROS~2, as relações espaciais entre referenciais são gerenciadas pela biblioteca
TF2. Para navegação autônoma com SLAM, três referenciais são centrais:

\begin{description}
    \item[\texttt{map}] Referencial global e estacionário, publicado pelo SLAM.
    A origem é tipicamente a posição do robô no instante de inicialização.

    \item[\texttt{odom}] Referencial de odometria, que acumula o deslocamento
    do robô via encoders desde a inicialização. É contínuo mas deriva.

    \item[\texttt{base\_\allowbreak link}] Referencial do robô propriamente dito,
    ancorado ao seu centro geométrico.
\end{description}

A transformação \texttt{map} $\to$ \texttt{odom} é publicada pelo SLAM e
representa a correção acumulada sobre a odometria pura. A transformação
\texttt{odom} $\to$ \texttt{base\_\allowbreak link} é publicada pelo driver de odometria
do robô. Juntas, formam a cadeia completa \texttt{map} $\to$ \texttt{odom}
$\to$ \texttt{base\_\allowbreak link}, que fornece a estimativa de pose do robô no
referencial do mapa com a maior precisão disponível.

\section{O Ecossistema ROS~2}
\label{sec:ros2}

O ROS~2 (\textit{Robot Operating System 2}) é uma plataforma de middleware para
robótica publicada como software livre pela Open Robotics \cite{macenski2022ros2}.
Desenvolvida desde 2014 para corrigir limitações arquiteturais do ROS~1, foi
construída sobre o protocolo DDS (\textit{Data Distribution Service}), que lhe
confere comunicação distribuída nativa, descoberta automática de nós, e suporte
a múltiplos perfis de qualidade de serviço (QoS).

\subsection{Políticas de QoS}
\label{sec:qos}

A configuração de QoS determina o comportamento da comunicação entre publicadores
e subscritores. Para experimentos de coleta de dados, três dimensões são críticas:

\begin{description}
    \item[\texttt{Reliability}] \textbf{BEST\_EFFORT} descarta mensagens se o
    buffer encher; \textbf{RELIABLE} garante entrega com retransmissão. Sensores
    de alta frequência (LiDAR, IMU) usam BEST\_EFFORT para minimizar latência.

    \item[\texttt{Durability}] \textbf{VOLATILE} descarta mensagens para subscritores
    que chegam após a publicação; \textbf{TRANSIENT\_LOCAL} mantém a última mensagem
    para novos subscritores. O \texttt{/\allowbreak amcl\_\allowbreak pose} usa TRANSIENT\_LOCAL para que
    um novo nó sempre receba a última estimativa de pose.

    \item[\texttt{History}] \textbf{KEEP\_LAST(N)} mantém as N mensagens mais
    recentes; \textbf{KEEP\_ALL} mantém todas. Tópicos de navegação crítica usam
    KEEP\_LAST(1) para garantir que apenas a informação mais recente seja processada.
\end{description}

A incompatibilidade de QoS entre publicador e subscritor é \textit{silenciosa}
no ROS~2: o match não é estabelecido, nenhuma mensagem é entregue, e nenhum
erro é reportado. Esta característica é responsável por uma categoria inteira
de bugs difíceis de depurar em experimentos de coleta de dados \cite{macenski2022ros2}.

\subsection{Modelo de Ciclo de Vida (Lifecycle)}
\label{sec:lifecycle}

O ROS~2 introduz um modelo de lifecycle para nós que executam inicialização
complexa \cite{macenski2022ros2}. Um nó lifecycle transita entre quatro estados:
\textit{unconfigured}, \textit{inactive}, \textit{active} e \textit{finalized}.
A transição \textit{unconfigured} $\to$ \textit{inactive} (\textit{configure})
carrega parâmetros e aloca recursos; a transição \textit{inactive} $\to$
\textit{active} (\textit{activate}) inicia o processamento.

O Nav2 usa extensivamente o lifecycle por meio do \texttt{lifecycle\_\allowbreak manager},
que ativa os servidores de planeamento, controle e comportamento em sequência,
verificando a saúde de cada um antes de prosseguir. Esse comportamento implica
que uma chamada prematura a um action server de navegação --- antes da ativação
completa --- resulta em timeout silencioso. O framework proposto nesta dissertação
implementa espera explícita pelo sinal de ativação completa antes de qualquer
operação de navegação.

\section{Navigation2 (Nav2)}
\label{sec:nav2}

O Nav2 é a pilha de navegação padrão para ROS~2 \cite{macenski2020marathon}.
Sua arquitetura decompõe o problema de navegação em servidores especializados
que se comunicam por meio de ROS~2 actions e serviços:

\begin{description}
    \item[\texttt{planner\_\allowbreak server}] Computa planos globais (trajetória completa
    até o destino) usando algoritmos como NavFn (Dijkstra/A*) ou Theta*.

    \item[\texttt{controller\_\allowbreak server}] Executa o plano localmente, publicando
    comandos de velocidade (\texttt{cmd\_\allowbreak vel}) para o robô. Na versão Jazzy,
    o controlador padrão é o MPPI (\textit{Model Predictive Path Integral})
    \cite{williams2016}, que amostra trajetórias candidatas e seleciona a que
    minimiza um funcional de custo. O Nav2 também suporta o Regulated Pure Pursuit
    \cite{macenski2023survey} e o Dynamic Window Approach \cite{fox1997} como
    controladores alternativos.

    \item[\texttt{behavior\_\allowbreak server}] Gerencia comportamentos de recuperação
    (\textit{spin}, \textit{backup}, \textit{drive on heading}) ativados quando
    o planeamento ou controle falham.

    \item[\texttt{bt\_\allowbreak navigator}] Orquestra os servidores via \textit{Behavior Trees}
    (BTs) \cite{colledanchise2018}. Cada ação de navegação corresponde a uma
    árvore de comportamento que define a sequência de operações, condições de
    falha e estratégias de recuperação.
\end{description}

\subsection{O Controlador MPPI}
\label{sec:mppi}

O MPPI (\textit{Model Predictive Path Integral Control}) \cite{williams2016} é um
algoritmo de controle ótimo baseado em amostragem. Em cada ciclo de controle, o
controlador gera $K$ trajetórias candidatas amostrando ruído gaussiano sobre a
sequência de controles atual; avalia o custo de cada trajetória segundo um funcional
que penaliza colisões, desvios do plano global e velocidades extremas; e atualiza
o controle com a média ponderada das trajetórias amostradas, com pesos inversamente
proporcionais ao custo.

Essa abordagem tem duas vantagens relevantes para experimentos de repetibilidade:
(i) produz trajetórias suaves e consistentes porque a atualização é um média ponderada,
não uma seleção discreta; (ii) é naturalmente paralelizável em GPU, embora em
aplicações com TurtleBot4 em CPU a versão CPU seja suficiente.

\subsection{Actions de Navegação}
\label{sec:actions}

O Nav2 expõe sua funcionalidade via ROS~2 actions --- um protocolo assíncrono
que permite enviar uma meta, receber feedback periódico e cancelar execuções em
andamento. Os dois actions centrais para este trabalho são:

\begin{itemize}
    \item \texttt{NavigateToPose}: navega o robô para uma única pose no mapa.
    Utilizado no modo de exploração manual e no retorno ao ponto de partida.

    \item \texttt{FollowWaypoints}: percorre sequencialmente uma lista de poses.
    Utilizado nas execuções de replay, garantindo que o robô visite exatamente
    os pontos gravados no baseline de forma contínua e coordenada.
\end{itemize}

A diferença comportamental entre os dois actions é relevante: o \texttt{NavigateToPose}
chamado sequencialmente para cada waypoint faz o robô parar completamente em cada
ponto antes de replanejar para o próximo, enquanto o \texttt{FollowWaypoints} encadeia
os planos, produzindo trajetórias mais suaves mas geometricamente diferentes. Essa
distinção entre baseline e replay constitui uma fonte de variação sistemática que
o framework identifica e documenta.


\subsection{Ciclo de Vida de uma Meta de Action}
\label{sec:action_lifecycle}

O protocolo de \textit{actions} do ROS~2 define uma máquina de estados bem
definida para cada meta enviada a um servidor: uma meta é primeiro aceita ou
rejeitada (\texttt{ACCEPTED}/\texttt{REJECTED}); em seguida, se aceita, evolui
para \texttt{EXECUTING}, podendo ser cancelada (\texttt{CANCELING} $\to$
\texttt{CANCELED}) ou concluir em \texttt{SUCCEEDED} ou \texttt{ABORTED}. Essa
máquina de estados é exposta ao cliente por meio de três \textit{futures}
assíncronos encadeados: o primeiro resolve quando o servidor aceita ou rejeita
a meta; o segundo, obtido a partir do \textit{goal handle} retornado pelo
primeiro, resolve quando a execução termina, com o \texttt{GoalStatus} final
anexado ao resultado.

Essa assincronia de três estágios tem uma implicação direta de implementação,
discutida em detalhe na Seção~\ref{sec:impl_orquestrador}: como o orquestrador
usa o mesmo mecanismo de \textit{callback} tanto para atualizar o estado interno
do robô (\texttt{nav\_\allowbreak state}) quanto para decidir se uma pose deve ser salva
durante uma gravação, um erro na composição desses \textit{callbacks} --- por
exemplo, tratar \texttt{ACCEPTED} como se fosse \texttt{SUCCEEDED} --- resultaria
em poses gravadas antes que o robô de fato tivesse alcançado o destino. A
implementação evita esse erro condicionando a gravação exclusivamente ao
\textit{callback} de resultado (Código~\ref{lst:start_stop_record} não mostra
esse trecho especificamente, mas o método \texttt{\_\allowbreak make\_\allowbreak result\_\allowbreak cb} descrito
na Seção~\ref{sec:impl_orquestrador} o implementa), nunca ao de aceitação.

\section{Localização: Filtros Bayesianos e AMCL}
\label{sec:localizacao_bayes}

Enquanto o SLAM Toolbox (Seção~\ref{sec:slam_toolbox}) resolve o problema de
\emph{construir} o mapa e estimar a pose simultaneamente, uma parcela relevante
dos sistemas de navegação autônoma opera com um mapa já conhecido e resolve
apenas o subproblema de localização --- estimar a pose do robô dado um mapa
fixo e um fluxo de observações sensoriais ruidosas. Esse é o papel do AMCL
(\textit{Adaptive Monte Carlo Localization}), o localizador padrão empacotado
com o Nav2 e utilizado como alternativa ao SLAM ao vivo nas análises de
sensibilidade do Capítulo~\ref{chap:resultados}.

O AMCL é uma instância de filtro de partículas: mantém um conjunto de $M$
hipóteses de pose (partículas), cada uma ponderada por sua verossimilhança
face à observação de \textit{laser scan} mais recente, e resample as partículas
proporcionalmente a esse peso a cada iteração. Formalmente, o filtro de
partículas aproxima numericamente a mesma recursão bayesiana que o filtro de
Kalman resolve de forma fechada sob hipóteses de linearidade e ruído gaussiano
\cite{kalman1960}: a etapa de predição aplica o modelo de movimento do robô a
cada partícula, e a etapa de atualização reponderada as partículas pela
verossimilhança da observação. A vantagem do filtro de partículas sobre o
filtro de Kalman clássico é justamente dispensar a hipótese de linearidade do
modelo de observação --- o modelo de \textit{laser scan matching} contra um
mapa de ocupação é fortemente não linear e multimodal (múltiplas poses podem
explicar igualmente bem a mesma leitura de LiDAR em ambientes com simetria,
como corredores repetidos de um armazém) --- às custas de um custo
computacional proporcional a $M$, tipicamente na ordem de centenas a milhares
de partículas.

Essa distinção tem relevância direta para o framework proposto, ainda que por um
motivo mais sutil do que uma leitura superficial da literatura sugeriria.
\citeonline{hamer2025} (discutida na Seção~\ref{sec:rel_es}) comparam o custo
computacional de quatro algoritmos de SLAM entre si --- não SLAM contra AMCL --- e
mostram que o paradigma grafo-vs-filtro não é, isoladamente, um bom preditor de
custo: o Cartographer, baseado em grafos, foi o mais caro em CPU (cerca de 43\% em
um Raspberry Pi~4) por fundir dados de IMU e realizar refinamento de mapa e
fechamento de laços, enquanto o Karto, também baseado em grafos, foi o algoritmo
mais eficiente do estudo em todas as métricas avaliadas. O fator determinante
identificado pelos autores é o conjunto de escolhas de projeto de cada
implementação --- fusão sensorial, frequência de publicação de mensagens,
refinamento de mapa --- não a família algorítmica em si.

Mesmo assim, o argumento estrutural desta dissertação para preferir AMCL com mapa
pré-construído em sessões experimentais longas permanece válido por um motivo mais
elementar, independente do algoritmo específico de SLAM utilizado: SLAM precisa
manter e expandir continuamente uma representação do ambiente (grafo de poses,
submapas, fechamento de laços) enquanto localiza o robô, ao passo que o AMCL
localiza um robô já posicionado sobre um mapa finalizado, com um conjunto de
partículas de tamanho fixo. Essa diferença estrutural --- construir e otimizar um
mapa versus apenas localizar-se nele --- é uma das razões que motivam a proposta de
usar AMCL com mapa pré-construído como alternativa ao SLAM ao vivo em sessões
experimentais longas, discutida como trabalho futuro na Seção~\ref{sec:futuros} a
partir da variabilidade observada nos experimentos exploratórios da
Seção~\ref{sec:res_exp_complementares}.

\section{SLAM Toolbox}
\label{sec:slam_toolbox}

O SLAM Toolbox \cite{macenski2021slam} é a implementação de SLAM 2D recomendada
para ROS~2. Trabalhos anteriores demonstraram que abordagens baseadas em grafos
superam filtros de partículas em ambientes de médio a grande porte
\cite{bailey2006}. Distingue-se de alternativas como o GMapping
\cite{grisetti2007} por três características:

\begin{enumerate}
    \item \textbf{Grafos de poses}: representa a trajetória do robô como um grafo
    e otimiza globalmente via fator Ceres, produzindo mapas de alta consistência
    mesmo em ambientes com fechamento de loop.

    \item \textbf{Mapeamento lifelong}: suporta sessões de mapeamento múltiplas,
    fundindo novos scans com mapas previamente construídos. Relevante para
    experimentos que reusam o mesmo ambiente ao longo de dias.

    \item \textbf{Modos sync/async}: o modo síncrono (\textit{online sync})
    processa cada scan antes de prosseguir, garantindo consistência absoluta
    entre scans e mapa. O modo assíncrono prioriza a taxa de publicação,
    adequado para robôs com CPU limitada.
\end{enumerate}

Para os experimentos desta dissertação, o modo síncrono é adotado para maximizar
a qualidade do mapa SLAM durante os experimentos de validação.

\section{Sensores e Coleta de Dados}
\label{sec:sensores}

O TurtleBot4 Standard integra três fontes de dados sensoriais relevantes para
análise de repetibilidade de navegação:

\subsection{LiDAR 2D (RPLIDAR S2)}

O RPLIDAR S2 é um sensor de varredura rotativo que emite pulsos laser em 360°
e mede distâncias por ToF (\textit{Time-of-Flight}). Opera a frequências de até
32\,kHz de amostragem e 10\,Hz de rotação. Cada scan resulta em uma mensagem
\texttt{sensor\_\allowbreak msgs/\allowbreak msg/\allowbreak LaserScan} com os campos \texttt{ranges} (distâncias em
metros por ângulo), \texttt{range\_\allowbreak min} e \texttt{range\_\allowbreak max} (limites de alcance
válido). O SLAM consome essas mensagens para atualizar o mapa e a estimativa de pose.

\subsection{Odometria de Roda}

A odometria é publicada como \texttt{nav\_\allowbreak msgs/\allowbreak msg/\allowbreak Odometry}, contendo a pose
estimada e a velocidade do robô. Em simulação Gazebo Harmonic, o plugin DiffDrive
computa a odometria a partir das velocidades das juntas e publica diretamente no
tópico \texttt{/\allowbreak odom} via bridge ROS-Gazebo. A frequência típica é 50\,Hz.

\subsection{IMU}

A IMU (\textit{Inertial Measurement Unit}) do Create3 publica aceleração linear
e velocidade angular a até 200\,Hz como \texttt{sensor\_\allowbreak msgs/\allowbreak msg/\allowbreak Imu}. Em experimentos
de repetibilidade, a variância da aceleração durante o percurso serve como indicador
indireto de instabilidade --- vibração acima do esperado pode indicar superfícies
irregulares ou acelerações abruptas do controlador \cite{siegwart2011}.

\section{Métricas de Repetibilidade}
\label{sec:metricas_teoricas}

A avaliação quantitativa de repetibilidade de trajetórias requer a comparação de
curvas espaço-temporais potencialmente de durações diferentes. Dois desafios
metodológicos principais emergem: (i) como amostrar as trajetórias de forma
comparável, e (ii) qual métrica reflete melhor o conceito intuitivo de ``mesma
trajetória''.

\subsection{Interpolação para Grade de Tempo Comum}
\label{sec:interpolacao}

Execuções de replay com o Nav2 raramente têm exatamente a mesma duração: variações
de até 5\% são comuns devido ao jitter do sistema operacional, variações de carga
computacional e o comportamento adaptativo do controlador MPPI. Para comparar
trajetórias de durações ligeiramente diferentes, é necessário normalizá-las para
uma grade de tempo comum \cite{amigoni2010}.

A abordagem adotada nesta dissertação normaliza cada trajetória para o intervalo
$[0, 1]$ e interpola linearmente os valores de $x$ e $y$ em $n = 100$ pontos
igualmente espaçados. A interpolação linear é adequada porque a frequência de
amostragem da odometria (50\,Hz) é muito maior que a dinâmica do robô, e entre
amostras consecutivas o movimento é essencialmente linear.

\subsection{RMSE Pairwise}
\label{sec:rmse_pairwise}

O RMSE (\textit{Root Mean Square Error}) entre duas trajetórias interpoladas
$T_i = \{(x_i^k, y_i^k)\}_{k=1}^{n}$ e $T_j = \{(x_j^k, y_j^k)\}_{k=1}^{n}$
é calculado como:

\begin{equation}
\text{RMSE}(T_i, T_j) = \sqrt{\frac{1}{n} \sum_{k=1}^{n}
\left[(x_i^k - x_j^k)^2 + (y_i^k - y_j^k)^2 \right]}
\label{eq:rmse_pairwise}
\end{equation}

Para uma frota com $N$ execuções, calcula-se a matriz $N \times N$ de RMSE
pairwise, onde o elemento $(i, j)$ representa a distância entre as execuções
$i$ e $j$. A diagonal é zero por construção; a simetria decorre da definição.

\subsection{Erro de Endpoint}
\label{sec:erro_endpoint}

O erro de endpoint complementa o RMSE com uma medida de onde o robô termina
o percurso:

\begin{equation}
e_{\text{end}}(T_i, T_j) = \sqrt{(x_i^n - x_j^n)^2 + (y_i^n - y_j^n)^2}
\label{eq:endpoint}
\end{equation}

Um erro de endpoint pequeno com RMSE grande indica que os percursos chegam ao
mesmo destino por caminhos diferentes --- comportamento aceitável para muitas
aplicações de navegação. A relação entre as duas métricas caracteriza o
tipo de variabilidade presente no sistema.

\subsection{Consistência Temporal}
\label{sec:consistencia_temporal}

A consistência temporal mede se execuções equivalentes têm durações similares.
Definimos o coeficiente de variação temporal entre $N$ execuções como:

\begin{equation}
\text{CV}_t = \frac{\sigma_t}{\bar{t}} \times 100\%
\label{eq:cv_temporal}
\end{equation}

\noindent onde $\bar{t}$ é a duração média e $\sigma_t$ o desvio padrão das
durações. Um $\text{CV}_t < 5\%$ indica consistência temporal suficiente para
que a interpolação para grade de tempo comum não introduza artefatos geométricos
significativos.

\section{Considerações Finais}
\label{sec:fund_conclusao}

Este capítulo estabeleceu os fundamentos teóricos que sustentam a proposta desta
dissertação, progressivamente da cinemática do robô às métricas de repetibilidade
de trajetórias. A transição entre esses extremos --- de encoders de roda ao RMSE
pairwise --- não é trivial; cada camada introduz abstrações, aproximações e fontes
de erro que o framework proposto deve tratar explicitamente.

O próximo capítulo examina como trabalhos existentes abordam partes desse problema,
e o que permanece em aberto após uma revisão cuidadosa da literatura.
